\section*{第 \ref{ch:g11:lithosphere-asthenosphere} 章——岩石圈与软流圈}

\begin{solution}{exo:g11:lithosphere-asthenosphere:1}
\omterm{def:g11:lithosphere-asthenosphere:lithosphere}{岩石圈}是地球坚硬的外壳，由地壳和最上部的
地幔组成；\omterm{def:g11:lithosphere-asthenosphere:lithosphere}{软流圈}是位于其下方的、温度更高、强度较弱的地幔，
它在漫长的时间尺度上缓慢流动。大洋地壳、
大陆地壳和最上部的地幔属于\omterm{def:g11:lithosphere-asthenosphere:lithosphere}{岩石圈}，外核
不属于\omterm{def:g11:lithosphere-asthenosphere:lithosphere}{岩石圈}。
\end{solution}

\begin{solution}{exo:g11:lithosphere-asthenosphere:2}
冰山下沉，直到排开的水重
等于自身重量：水下部分所占比例为
$910/1025 \approx 0.888$，约 \qty{89}{\percent}。露出水面部分约占
\qty{11.2}{\percent}，因此总高度为 $20/0.112 \approx \qty{180}{m}$，
其中约 \qty{160}{m} 位于水下。
\end{solution}

\begin{solution}{exo:g11:lithosphere-asthenosphere:3}
$15 + 26 \times 3 = \qty{93}{\celsius}$。在 \qty{100}{km} 处，若仍采用相同梯度，
温度将超过 \qty{2600}{\celsius}，远高于地幔岩石的熔化温度；但 \omterm{def:g11:seismic-structure:body-waves}{S 波}能穿过上地幔，
可见上地幔是固态的。因此，
\omterm{def:g11:lithosphere-asthenosphere:geotherm}{地温梯度}必须随深度显著减小。
\end{solution}

\begin{solution}{exo:g11:lithosphere-asthenosphere:4}
海平面以上 0 至 \qty{500}{m} 的区间，约占地表的 \qty{16}{\percent}，
以及海平面以下 4000 至 \qty{4500}{m} 的区间，约占
\qty{11}{\percent}（海平面以下 4500 至 \qty{5000}{m}
的区间几乎同样大）。
\end{solution}

\begin{solution}{exo:g11:lithosphere-asthenosphere:5}
当 $P = 0$ 时：\qty{1121}{\celsius}。当 \qty{3}{GPa} 时：$-5.104 \times 9 +
132.899 \times 3 + 1120.661 \approx -45.9 + 398.7 + 1120.7 \approx
\qty{1473}{\celsius}$。
\end{solution}

\begin{solution}{exo:g11:lithosphere-asthenosphere:6}
在 \qty{50}{km} 深处约为 \qty{730}{\celsius}（\omterm{def:g11:lithosphere-asthenosphere:lithosphere}{岩石圈}），
在 \qty{150}{km} 深处约为 \qty{1370}{\celsius}，在 \qty{250}{km}
深处为 \qty{1410}{\celsius}（\omterm{def:g11:lithosphere-asthenosphere:lithosphere}{软流圈}）。温度升高最快的是
50 至 \qty{150}{km} 之间，升高了 \qty{600}{\celsius} 以上，因为这段区间
大部分位于\omterm{def:g11:lithosphere-asthenosphere:lithosphere}{岩石圈}内；在 150 至 \qty{250}{km} 之间，温度仅约升高
\qty{45}{\celsius}。
\end{solution}

\begin{solution}{exo:g11:lithosphere-asthenosphere:7}
根部：$3 \times 2830/550 \approx \qty{15.4}{km}$。
增加的厚度：$3 \times 3380/550 \approx \qty{18.4}{km}$，
即山脉的 \qty{3}{km} 高度加上其山根。
\end{solution}

\begin{solution}{exo:g11:lithosphere-asthenosphere:8}
地壳增厚 $55 - 41.1 = \qty{13.9}{km}$，且 $\Delta H = h
\times 3380/550$，故 $h = 13.9 \times 550/3380 \approx \qty{2.3}{km}$。
\end{solution}

\begin{solution}{exo:g11:lithosphere-asthenosphere:9}
减小量：$(4.470 - 4.419)/4.470 \approx \qty{1.1}{\percent}$。跃升量：
$(4.644 - 4.419)/4.419 \approx \qty{5.1}{\percent}$。随深度增加，岩石受到的压缩越来越强，
地震波速随之增大；速度减小则说明另有因素起主导作用：岩石温度高，
接近熔化，强度较弱。
\end{solution}

\begin{solution}{exo:g11:lithosphere-asthenosphere:10}
(1) $(1345 - 610)/60 \approx \qty{12}{\celsius}$ 每千米。(2) 这约为
地表梯度的一半。部分热量产生于上地壳内部，这些热量同样需要穿过
顶部各层才能逸出，
因此近地表的梯度较陡。在地壳以下，穿过岩石的只是
来自更深处的热量，梯度因而较小。
\end{solution}

\begin{solution}{exo:g11:lithosphere-asthenosphere:11}
在 \qty{200}{km} 处：$1300 + 90 = \qty{1390}{\celsius}$；在 \qty{400}{km} 处，
$1300 + 180 = \qty{1480}{\celsius}$。可见自 \qty{200}{km} 至
\qty{400}{km}，温度仅升高 \qty{90}{\celsius}。穿过\omterm{def:g11:lithosphere-asthenosphere:lithosphere}{岩石圈}时，温度在
仅 \qty{100}{km} 内就升高 $1345 - 15 = \qty{1330}{\celsius}$，
每千米的升温速率约为前者的三十倍。
\end{solution}

\begin{solution}{exo:g11:lithosphere-asthenosphere:12}
(1) $118 + 12.3\,z = 1300 + 0.45\,z$ 可得 $z = 1182/11.85 \approx
\qty{100}{km}$。(2) $118 + 21.5\,z = 1300 + 0.45\,z$ 可得 $z = 1182/21.05
\approx \qty{56}{km}$。(3) 较热的西部\omterm{def:g11:lithosphere-asthenosphere:lithosphere}{岩石圈}明显更薄，
与实测相符（最薄处约 \qty{40}{km}）；大陆中部的\omterm{def:g11:lithosphere-asthenosphere:lithosphere}{岩石圈}
极厚，所需梯度还要更小。
\end{solution}

\begin{solution}{exo:g11:lithosphere-asthenosphere:13}
(1) 由命题可得 $\Delta H = h\,\rho_m/(\rho_m - \rho_c)$，故 $h =
\Delta H\,(\rho_m - \rho_c)/\rho_m$。(2) $550/3380 \approx 0.16$：在任何增加的地壳厚度中，
约六分之一表现为地表高度，其余六分之五表现为山根。(3) 对于
冰山，露出水面部分所占比例为 $(1025 - 910)/1025 \approx 0.11$：
同一规律为：（流体密度 $-$ 漂浮物体密度）/
流体密度。
\end{solution}

\begin{solution}{exo:g11:lithosphere-asthenosphere:14}
(1) $1489 - 1345 = \qty{144}{\celsius}$。(2) 少量熔体聚集在岩石颗粒之间
的薄膜中。颗粒之间不再紧密咬合，岩石抵抗剪切的能力减弱，
因此 \omterm{def:g11:seismic-structure:body-waves}{S 波}传播速度变慢，岩石更容易发生变形，
并能更轻易地流动。
\end{solution}

\begin{solution}{exo:g11:lithosphere-asthenosphere:15}
(1) 否。将陆地削平只会使大陆地壳稍微变薄，
但它仍然厚而轻，
仍然比薄而致密的洋壳浮得高得多；
海洋依然会充满大陆之间的海盆。(2) 地壳分为两类，
厚度与密度各不相同：厚而轻的大陆地壳，
以及薄而致密的洋壳。
\end{solution}

\begin{solution}{pb:g11:lithosphere-asthenosphere:1}
\textbf{1.} $5.01 - 0.67 = \qty{4.34}{km}$.

\textbf{2.} $(3380 - 2830)/3380 \approx \qty{16}{\percent}$.

\textbf{3.} 若按分数 $2830/3380 \approx 0.84$ 计算，冰山应有此比例的部分位于水面以下；
实际上冰山沉得稍深，$910/1025 \approx 0.89$，因为冰与水的
密度差比地壳与地幔之间的密度差更小。

\textbf{4.} 在地震波传播的秒级时间内，或在一次地震的时间内，
地幔表现为固体；只有在极长的时间尺度上，
\omterm{def:g11:lithosphere-asthenosphere:lithosphere}{软流圈}才会流动，使地壳得以漂浮。

\textbf{5.} 在最深的\omterm{def:g11:seismic-structure:moho}{莫霍面}以下，两个柱体均由相同的地幔岩石构成，
在相同深度处重量相同；所有差异
都位于该深度以上，必须计入。

\textbf{6.} 平均大陆：地壳从海平面以上 \qty{0.67}{km} 延伸至
海平面以下 \qty{40.4}{km}，其下为地幔，直至 \qty{70}{km}。
高原：地壳自海平面以上 \qty{5.01}{km} 延伸至
因山根而加深的\omterm{def:g11:seismic-structure:moho}{莫霍面}，其下为地幔，直至 \qty{70}{km}。

\textbf{7.} 高原柱体多出 $h + r$ 的地壳，大陆柱体
则在厚度 $r$ 内为地幔：$\rho_c\,(h + r) = \rho_m\,r$，
于是 $r = h\,\rho_c/(\rho_m - \rho_c)$，且 $\Delta H = h + r =
h\,\rho_m/(\rho_m - \rho_c)$。

\textbf{8.} $3380/550 \approx 6.15$.

\textbf{9.} $\Delta H = 4.34 \times 6.15 \approx \qty{26.7}{km}$；山根厚度为
$26.7 - 4.34 \approx \qty{22.3}{km}$。

\textbf{10.} 平均大陆的\omterm{def:g11:seismic-structure:moho}{莫霍面}位于海平面以下 $41.1 - 0.67 \approx
\qty{40.4}{km}$；高原的\omterm{def:g11:seismic-structure:moho}{莫霍面}位于海平面以下 $40.4 + 22.3 \approx
\qty{62.8}{km}$。

\textbf{11.} $41.1 + 26.7 \approx \qty{67.8}{km}$.

\textbf{12.} 每平方米之上一千米厚的岩石，质量为
$\rho \times \qty{1000}{m^3}$：地壳 \num{2830} 吨，或地幔 \num{3380}
吨。大陆柱体：$41.1 \times 2830 + (70 - 40.43) \times 3380
\approx \num{116300} + \num{99900} \approx \num{216300}$ 吨。高原柱体：
$67.77 \times 2830 + (70 - 62.76) \times 3380 \approx \num{191800} +
\num{24500} \approx \num{216300}$ 吨。两者相等。

\textbf{13.} $41.1 - 0.67 \times 6.15 \approx \qty{37.0}{km}$.

\textbf{14.} $67.8/41.1 \approx 1.65$ 倍。

\textbf{15.} $3300/470 \approx 7.02$, $\Delta H \approx 4.34 \times 7.02
\approx \qty{30.5}{km}$，地壳约为 \qty{71.6}{km}。地幔密度的
变化约 \qty{2}{\percent}，会使厚度改变
约 \qty{4}{km}：结果明显依赖于密度，但仍
接近 \qty{70}{km}。

\textbf{16.} 模型给出约 \qty{68}{km}，地震波给出约
\qty{70}{km}：两者在百分之几以内相符。

\textbf{17.} $h = 3.23 - 0.67 = \qty{2.56}{km}$, $\Delta H = 2.56 \times
6.15 \approx \qty{15.7}{km}$，地壳约为 \qty{56.8}{km}，比实测值
多约 \qty{7}{km}。

\textbf{18.} 更薄。柱体上部比平均值轻，
在厚度相同时重量更小，因此浮得更高：较轻的地壳
要达到给定高度所需山根更小。实际地壳因此可以
比模型所给的更薄，正如实测所示。

\textbf{19.} 在长时间尺度上，\omterm{def:g11:lithosphere-asthenosphere:lithosphere}{软流圈}会流动：它在
较重柱体之下向旁侧移开，又回流到较轻柱体之下，直至两柱达到平衡。

\textbf{20.} \textbf{西藏之下的地壳厚约 \qty{70}{km}}，远大于
平均大陆的 \qty{41}{km}：高原高出海平面约五千米，
因为其厚而轻的地壳漂浮在密度更大的地幔之上，
山根深度超过 \qty{20}{km}。
\end{solution}
